\documentclass[11pt,a4paper]{article}
\usepackage[margin=27mm]{geometry}
\usepackage{amsmath,amssymb,amsthm,fontspec,unicode-math,hyperref,booktabs,longtable}
\setmainfont{STIXGeneral}
\setmathfont{STIX Two Math}
\setmonofont{Menlo}
\usepackage{newunicodechar}
\newunicodechar{∎}{\ensuremath{\blacksquare}}
\hypersetup{hidelinks}
\setlength{\emergencystretch}{2em}
\providecommand{\tightlist}{\setlength{\itemsep}{0pt}\setlength{\parskip}{0pt}}
\newcommand{\gen}[1]{\langle #1\rangle}
\newcommand{\Sym}{\operatorname{Sym}}
\newcommand{\Alt}{\operatorname{Alt}}
\newcommand{\core}{\operatorname{core}}
\newcommand{\F}{\mathbb F}
\newtheorem{theorem}{Theorem}
\newtheorem{lemma}[theorem]{Lemma}
\newtheorem{proposition}[theorem]{Proposition}
\newtheorem{corollary}[theorem]{Corollary}

\title{Frattini quotients and primitive coordinate geometries}
\author{Human Compression project}
\date{}
\begin{document}
\maketitle
\begin{abstract}
Firm residually weakly primitive geometries of finite rank factor through the Frattini quotient, even when the acting group is infinite. For finite $p$-groups the complete subgroup family is the pullback of coordinate hyperplanes, giving the classification, its converse and the precise count of indexed families. We prove the relative-generation and quotient interfaces with their exact hypotheses and retain an explicit counterexample to the unrestricted infinite generating-set assertion.
\end{abstract}
\section{Introduction}
Archer, Cara and Krempa used the Frattini subgroup to reduce the study of residually weakly primitive geometries. The geometric containment theorem and the relative generating-set theorem have different hypotheses in infinite groups. We keep that distinction explicit. The common paper [UB] proves flag transitivity and exact relative parabolic generation once; the complete normal-subgroup ladder, quotient correspondence and coordinate adapters needed here are supplied in this paper. The field is $\mathbb F_p$ for an arbitrary prime $p$, not only for $p=2$.
\section{Relative generation and removal of normal kernels}
\textbf{Proposition B3.Quotient (exact coset transport).} Let \((H_i)_{i\in I}\) be subgroups of \(G\), let \(B=\bigcap_i H_i\), and let \(N\unlhd G\) satisfy \(N\leq B\). Quotienting by \(N\) gives a bijection \begin{equation}
H_i g\longmapsto (H_i/N)(gN)
\label{eq:frattini-coordinate-geometries:1-5}
\end{equation} on the typed right cosets. Taking inverses gives exactly the left-coset convention of [UB]. It preserves the intersections of arbitrary families of those cosets, including the existence of a common element. It consequently preserves incidence, flags, residues and all properties determined by the resulting incidence system. The actions of \(G\) and \(G/N\) have the same permutation image. Maximality of any inclusion between intersections of the \(H_i\) is preserved.

\emph{Proof.} The inverse image of the coset on the right of \eqref{eq:frattini-coordinate-geometries:1-5} is exactly \(H_i g\). Thus two source cosets give the same target coset precisely when they were equal, and taking inverse images commutes with all intersections. Surjectivity of the quotient map supplies a source representative of each target coset and a common lift whenever a target intersection is nonempty. This proves the assertions about incidence and common elements. The quotient action is visibly the same action on the identified coset set. Every subgroup in any interval under consideration contains \(N\); the subgroup correspondence therefore preserves that interval and its maximal inclusions. \(\square\)
\section{The normal Frattini ladder}
\textbf{Lemma B3.Maximal (relative finite generation).} If \(G=\langle B,t_1,\ldots,t_r\rangle\) with \(r\) finite, then every proper subgroup containing \(B\) is contained in a maximal subgroup of \(G\).

\emph{Proof.} Order the proper overgroups of the given subgroup by inclusion. The union of a chain is a subgroup. If that union were \(G\), a single member of the chain would contain all the finitely many \(t_j\), and hence all of \(G\), a contradiction. Thus every chain has a proper upper bound. Zorn's lemma gives a maximal proper overgroup. \(\square\)

Define \(\Phi(G)\) to be the intersection of all maximal subgroups, with intersection \(G\) if there are none. This subgroup is characteristic. Whenever every proper overgroup of \(B\) is contained in a maximal subgroup, a subgroup \(H\geq B\) satisfies \begin{equation}
\Phi(G)H=G\quad\Longleftrightarrow\quad H=G:
\label{eq:nongeneration}
\end{equation} a maximal overgroup of a proper \(H\) would contain both factors on the left.

\textbf{Theorem B3.Transport (relative generating families).} Under the maximal-overgroup hypothesis of \eqref{eq:nongeneration}, a finite indexed family is an irredundant relative generating family of \(G\) over \(B\) if and only if its image is one of \(G/\Phi(G)\) over \(B\Phi(G)/\Phi(G)\).

\emph{Proof.} Generation in the quotient says \(\Phi(G)\langle B,s_i:i\in I\rangle=G\), which is equivalent to generation in \(G\) by \eqref{eq:nongeneration}. Once generation holds, irredundancy says that each deletion subgroup \(\langle B,s_j:j\ne i\rangle\) is proper. Equation \eqref{eq:nongeneration}, applied to each of those subgroups, says exactly that its image is proper in the quotient. These are precisely the quotient irredundancy conditions. \(\square\)

The maximal-overgroup hypothesis holds for finite \(G\) and, by B3.Maximal, for relative finite generation. It cannot simply be omitted for arbitrary infinite groups: in \(C_{p^\infty}\times C_2\), the divisible factor is the Frattini subgroup, and a generator of the quotient does not generate the whole group.

\textbf{Theorem B3.Ladder (containment in the common stabilizer).} Suppose the \(H_J\) are all distinct and some ordering of \(I=\{1,\ldots,r\}\) gives a chain \begin{equation}
G=K_0>K_1>\cdots>K_r=B,\qquad K_j=H_{\{1,\ldots,j\}},
\label{eq:frattini-coordinate-geometries:6-6}
\end{equation} in which every inclusion is maximal. If \(N\unlhd G\) lies in every maximal subgroup of \(G\), then \(N\leq B\).

\emph{Proof.} Choose \(t_j\in K_{j-1}\setminus K_j\). Maximality gives \(K_{j-1}=\langle K_j,t_j\rangle\), and iteration gives \(G=\langle B,t_1,\ldots,t_r\rangle\). Every proper \(H_i\) therefore lies in a maximal subgroup by B3.Maximal. Such a maximal subgroup contains \(N\), so \(NH_i\ne G\).

Suppose \(N\not\leq B\), and choose the first \(k\) with \(N\not\leq K_k\). Then \(N\leq K_{k-1}\), and normality and maximality give \[
K_{k-1}=N(K_{k-1}\cap H_k).
\] Suppose at an upward step that \(K_j=N(K_j\cap H_k)\), where \(1\leq j<k\). Since \(N\leq K_{j-1}\), \[
K_j\leq N(K_{j-1}\cap H_k)\leq K_{j-1}.
\] If this middle group were \(K_j\), then \(K_{j-1}\cap H_k\leq K_j\), forcing \[
K_{j-1}\cap H_k=K_j\cap H_k.
\] The corresponding index sets differ because \(j<k\), contradicting the assumed distinctness. Maximality therefore makes the middle group \(K_{j-1}\). Continuing upward reaches \(G=NH_k\), contradicting the first paragraph. Hence \(N\leq B\). \(\square\)
\section{The complete geometry and coordinate classification}
The groups in this section may be infinite, with finite type set, except where a finite p-group is expressly specified. Put \(\Phi(G)=\bigcap\{M:M\text{ maximal in }G\}\), with an empty intersection equal to \(G\).

\textbf{Theorem B3.Frattini.} For a finite-rank firm RWPRI geometry with Borel \(B\), \begin{equation}
\Phi(G)\leq B\leq H_i\quad(i\in I).
\label{eq:frattini-coordinate-geometries:6-1}
\end{equation} The geometry is isomorphic, with the same permutation action, to the coset geometry of \(G/\Phi(G)\) and the subgroups \(H_i/\Phi(G)\). More generally this coset correspondence holds for any normal subgroup \(N\leq B\), even for a pregeometry without flag transitivity. The full action kernel is \begin{equation}
\operatorname{core}_G(B)=\bigcap_i\operatorname{core}_G(H_i).
\label{eq:frattini-coordinate-geometries:6-2}
\end{equation} Quotienting by it makes the action faithful.

\emph{Proof.} By UB.Flags and firmness all indexed intersections are distinct. Repeatedly choosing a maximal next intersection under (G4) supplies a chain from \(G\) to \(B\), every step maximal. The subgroup \(\Phi(G)\) is characteristic and lies in every maximal subgroup of \(G\). Thus all hypotheses of B3.Ladder hold, proving \eqref{eq:frattini-coordinate-geometries:6-1}, including for infinite \(G\). B3.Quotient now applies because the normal subgroup lies in every type stabilizer. Its coset bijection preserves all intersection patterns and maximal subgroup intervals, hence the geometry, residues and primitive properties.

An element acts trivially on every type precisely when it lies in the core of every \(H_i\). Intersecting all conjugates commutes with the finite intersection defining \(B\), proving \eqref{eq:frattini-coordinate-geometries:6-2}. This is the largest normal subgroup contained in \(B\), and B3.Quotient gives the faithful quotient. ∎

The quotient \(G/\Phi(G)\) is Frattini free: its maximal subgroups correspond to the maximal subgroups of \(G\), with intersection \(\Phi(G)\). If there are no maximal subgroups then the quotient is trivial, and the conclusion is the same. In particular every firm RWPRI geometry has a realization for a Frattini-free group. Pullback and quotient give a bijection of indexed subgroup families with the quotient map fixed. This does not assert that automorphisms of the quotient lift, or identify orbit counts under different automorphism groups.

For relative generating families, B3.Transport says that a finite family generates \(G\) irredundantly relative to \(B\) if and only if its image does so in \(G/\Phi(G)\), provided every proper overgroup of \(B\) is contained in a maximal subgroup. This hypothesis holds for finite groups. It also follows from B3.Maximal when \(G\) is finitely generated relative to \(B\). The unrestricted infinite converse printed in {[}ACK05, Theorem 2.1{]} requires this correction. Indeed let \(D=C_{p^\infty}\) and \(G=D\times C_2\). A maximal subgroup of this abelian group has a cyclic quotient of prime order. Divisibility makes the image of \(D\) in every such quotient trivial; the only maximal subgroup is therefore \(D\times1=\Phi(G)\). The generator of the second factor generates \(G/\Phi(G)\) but not \(G\). This counterexample does not restrict Theorem B3.Frattini: its maximal-overgroup argument is already furnished by the finite maximal chain in B3.Ladder.

\textbf{Theorem B3.Coordinates.} Let \(P\) be a finite p-group acting on a firm RWPRI geometry of rank \(r\), with Borel \(B\). Then \(B\unlhd P\), \(P/B\cong\mathbb F_p^r\), and there is a basis \((v_i)_{i\in I}\) such that \begin{equation}
H_J/B=\operatorname{span}\{v_i:i\notin J\}\quad(J\subseteq I).
\label{eq:frattini-coordinate-geometries:6-3}
\end{equation} The geometry is the direct sum of \(r\) primitive rank-one geometries of \(p\) points. It is RPRI, and its faithful group \(P/B\) is regular on chambers.

\emph{Proof.} B3.Frattini gives \(\Phi(P)\leq B\). The classical finite-$p$-group fact that \(P/\Phi(P)\) is elementary abelian identifies \(P/\Phi(P)\) with a vector space. Consequently \(B/\Phi(P)\) is normal, \(B\unlhd P\), and \(P/B\) is a vector space. Apply UB.Generation: the images \(v_i=s_iB\) span \(P/B\). A nonzero linear dependence expresses one \(s_i\) modulo \(B\) in terms of the other generators, contrary to relative independence. Thus they form a basis, and the parabolic formula of Theorem UB.Generation becomes \eqref{eq:frattini-coordinate-geometries:6-3}.

The cosets of \(H_i/B\) are exactly the \(p\) possible values of coordinate \(i\). Assignments to distinct coordinates always have a common vector. Hence every choice of one element of each type is a chamber, with all distinct types completely incident. Fixing one more coordinate in any residue has index \(p\), and is a maximal subgroup inclusion. This proves RPRI and the primitive direct-sum description. Since \(B\) is normal it is the action kernel in \eqref{eq:frattini-coordinate-geometries:6-2}; translations of \(P/B\) are regular on the set of complete coordinate assignments. ∎

Conversely every family of \(r\) independent functionals on \(P/\Phi(P)\) yields such a geometry by pulling back their kernels. The coordinate argument proves firmness, flag transitivity, residual connectedness and RPRI. Theorem B3.Coordinates shows there are no other indexed subgroup families. Writing \(d=\dim P/\Phi(P)\), their number is \begin{equation}
\frac{\prod_{j=0}^{r-1}(p^d-p^j)}{(p-1)^r},\qquad 1\leq r\leq d,
\label{eq:frattini-coordinate-geometries:6-4}
\end{equation} since one chooses independent ordered functionals and divides by the independent nonzero rescalings that preserve their kernels. This counts indexed subgroup families, not automorphism orbits.

Two subsidiary assertions of {[}ACK05{]} follow directly in these coordinates. For arbitrary subgroups \(H_1,H_2\), the cosets \(gH_1,hH_2\) meet precisely when \(g^{-1}h\in H_1H_2\). Thus their rank-two pregeometry is a direct sum exactly when \(H_1H_2=G\). If \(U,W<\mathbb F_p^k\), with \(U\) a hyperplane, \(U\cap W\) a hyperplane of \(U\), and \(W\not\subseteq U\), then \(W\) has one-dimensional image in \(\mathbb F_p^k/U\). It follows that \(\dim W=(k-2)+1=k-1\), as asserted in the source's maximal-subspace lemma.

Finally, every residue gives its own necessary Frattini test. At a flag of type \(J\), its stabilizer is \(H_J\); its type subgroups are \(H_J\cap H_i\), \(i\notin J\), and their common intersection is still \(B\). UB.Flags, firmness and residual weak primitivity hold within this residue. Applying B3.Ladder to these actual subgroups gives \begin{equation}
\Phi(H_J)\leq B\leq H_J\cap H_i\quad(i\notin J).
\label{eq:residue-frattini}
\end{equation} This is a necessary rejection test, not a sufficiency criterion. It does not rely on the generally false inclusion \(\Phi(H_J)\leq\Phi(G)\).
\section{The admitted source statements}
\textbf{Corollary B3.Source (primitive geometries and Frattini quotients).} A finite-rank firm RWPRI geometry, including one for an infinite group, has its Frattini subgroup inside every type stabilizer and is exactly the corresponding coset geometry of the Frattini quotient. Normal subgroups of the Borel can be removed without changing even the pregeometry; removing its core makes the action faithful. The quotient group is Frattini free. For a finite p-group, RWPRI is equivalent to RPRI, and the geometry is a direct sum of primitive rank-one geometries, with its complete coordinate classification given by B3.Coordinates. Every residue satisfies \eqref{eq:residue-frattini}. Relative irredundant generating families correspond through the Frattini quotient under B3.Transport's maximal-overgroup hypothesis, in particular for finite groups.

\emph{Proof.} These are B3.Frattini, B3.Coordinates and the residue application of B3.Ladder. They give {[}ACK05, Theorems 3.1, 3.4, 3.8, Proposition 3.2, Corollary 3.3, Proposition 3.7 and §§4.2--4.3{]}. The coset-product and subspace calculations following B3.Coordinates give Proposition 3.5 and Lemma 3.6. B3.Transport with the explicitly checked maximal-overgroup hypothesis gives Theorem 2.1 in the corrected scope; the paragraph following B3.Frattini gives the counterexample to its unrestricted infinite converse. This correction does not narrow the infinite geometric containment theorem. ∎

\section*{References and retained inputs}
{[}UB{]} \emph{Forest coordinates and relative generation}, common theory in the present volume. Its full proofs and definitions are part of this application's mathematical dependencies.

{[}ACK05{]} C. Archer, P. Cara and J. Krempa, \emph{Using the Frattini Subgroup and Independent Generating Sets to Study RWPRI Geometries}, Beiträge zur Algebra und Geometrie \textbf{46}(1) (2005), 169--177. The complete published nine-page PDF is the consulted source. Its elementary finite-p-group Frattini quotient fact is retained from p.~174. B3.Maximal uses Zorn's lemma in the infinite-group setting.

The structural theorem quoted as {[}ACK05, Theorem 4.1{]} and the older primitive-group reduction cited at the end of its §4.3 remain earlier work and lie outside the target clauses. Neither is an input to a proof above.

\end{document}
